\section{Platform and Environment Control}
\label{sec:platform}

\begin{figure}[t]
\centering
\begin{tikzpicture}[x=6.2mm,y=6mm,font=\scriptsize]
  \foreach \i/\c in {0/0,1/2,2/4,3/6,4/8,5/10,6/12,7/14} {
    \draw[fill=native!15, draw=native, thick] (\i*1.5,0) rectangle ++(1.2,1.2);
    \node at (\i*1.5+0.6,0.6) {\c};
  }
  \draw[fill=red!12, draw=native, thick] (9,0) rectangle ++(1.2,1.2);
  \node at (9.6,0.6) {12};
  \node[anchor=south] at (5.85,1.3) {8 P-cores, one CPU per core (SMT off): the measured mask};
  \foreach \i in {0,...,15} {
    \draw[fill=gray!15] (\i*0.75,-2.2) rectangle ++(0.6,0.9);
  }
  \node[anchor=north] at (6,-2.3) {16 E-cores (CPUs 16 to 31): the system, the harness, the results server};
  \draw[dashed] (-0.3,-0.35) -- (12.3,-0.35);
  \node[anchor=west] at (12.4,0.6) {run scope};
  \node[anchor=west] at (12.4,-1.75) {shield};
\end{tikzpicture}
\caption{The CPU layout during a measurement. A run is confined to the eight
P-cores by a systemd scope and a process-level mask; everything else is
confined to the E-cores. CPU 12 is the degraded core
(Section~\ref{sec:core}).}
\Description{Eight boxes for the P-core CPUs 0 to 14 and sixteen smaller boxes for the E-cores, separated by a dashed line.}
\label{fig:cores}
\end{figure}

\paragraph{Machine.} An Intel Core i9-14900KF on an ASUS ProArt Z790-CREATOR
WIFI (BIOS 3202), 64~GB of RAM, Ubuntu 26.04.1 LTS, kernel 7.0.0-43,
microcode 0x133. The CPU is hybrid: 8 P-cores with two hardware threads
each and 16 E-cores. The last-level cache is 36~MB.

\paragraph{P-cores only, one mask for all.} The study measures how the time
changes with $N$. If threads could land on either kind of core, the speed of
a thread would change with $N$ for reasons that have nothing to do with
WASM. The usual control, pinning each thread, is not available: a web page
cannot choose the core of a worker. So both sides get the same
process-level mask, the eight P-cores, and the kernel places threads inside
it (Figure~\ref{fig:cores}). A pilot confirmed the choice
(Section~\ref{sec:pilots}).

\paragraph{SMT off.} Linux numbers the two hardware threads of a P-core side
by side, so a mask such as \code{0-7} is four cores, not eight. SMT is off
for every session and the mask is \code{0,2,4,6,8,10,12,14}; every run
records the SMT state and the siblings of each masked CPU, and a run whose
mask has two CPUs on one core is rejected.

\paragraph{Turbo off, governor \code{performance}.} With turbo on, one busy
core runs at a higher clock than eight busy ones, so $T_1$ would be measured
at a higher frequency than $T_8$ and the speedup would be deflated for a
reason unrelated to parallelism. With turbo off the P-cores run at their
3.2~GHz base clock (the recorded maximum of every measured CPU is
3\,200\,000~kHz).

\paragraph{Shield.} A mask keeps the run on the P-cores; it does not keep
anything else off them. On a hybrid CPU the scheduler even prefers the
P-cores. The system, user and init slices are therefore confined to the
E-cores at run time, the harness and the results server pin themselves
there, and each run is started in its own systemd scope whose allowed CPUs
are the mask. What the shield cannot stop (per-CPU kernel threads,
interrupts) is measured: every run records the busy time of the measured
CPUs minus its own CPU time, the \emph{foreign CPU time}, and a run above a
threshold fixed from pilot data is excluded (Section~\ref{sec:pilots}).

\paragraph{Quiet mode.} Before a campaign the desktop session and bursty
services (containers, package timers, snaps) are stopped, because they share
the last-level cache and the memory bandwidth with the measured cores even
when they run on other CPUs; the machine does not sleep during a campaign.

\paragraph{Checked before every start, recorded with every run.} None of
these settings survives a reboot. A campaign refuses to start, and to
resume, unless turbo is off, the governor is \code{performance} on every
measured CPU, SMT is off, the mask is exactly one CPU per P-core, the shield
holds and the pinned browser versions are the installed ones. Each run
stores the commit, the CPU model, kernel, governor, turbo and SMT state, the
mask, the browser version, the launch command, the order seed and position,
the foreign CPU time, the processes found outside its scope, the load
average, and for every measured CPU its busy time. \emph{Strict validation}
rejects a run in which any of them is missing or wrong. Since v3 the host
state that the setup leaves alone is recorded too: transparent huge pages
(\code{madvise}), the kernel command line, the mitigations, the idle driver
and its C-states, the firmware and the memory size.

\subsection{The degraded core}
\label{sec:core}

One P-core of the machine is degraded: the seventh, CPU 12 of the mask.
Before the v3 build every P-core was checked, with the rule fixed in
advance: the eight cores stay whatever times the check shows, and a core
leaves the mask only on a wrong checksum or a crash.

\begin{itemize}[leftmargin=1.5em]
\item \emph{On the machine as found it fails.} With turbo on, SMT on, the
governor \code{powersave} and no shield, a test suite pinned to that core
crashed with a segmentation fault in 17 of 18 runs, and in 0 of 28 runs
pinned elsewhere.
\item \emph{On the prepared machine it does not.} After the setup the same
tests passed 30 of 30 times on CPU 12. The setup changes four things
together, so the check does not say which one matters; turbo is the likely
one, since it is what lets the core run above 3.2~GHz.
\item \emph{On the prepared machine it is as fast as the others}
(Table~\ref{tab:core-check}): the native filter at one thread pinned to one
core at a time gives 22.41 to 22.53~ms on every core, 22.44 on CPU 12, the
right checksum in all 80 runs, and no iteration above 22.9~ms among 800.
\end{itemize}

\begin{table}[htbp]
\caption{Per-core check on the prepared machine: the native filter at one
thread, 8192$\times$8192, pinned to one P-core at a time, 10 runs per core
in rotating order (a check, not a result).}
\label{tab:core-check}
\centering\small
\begin{tabular}{lrrrrrrrr}
\toprule
CPU & 0 & 2 & 4 & 6 & 8 & 10 & 12 & 14 \\
\midrule
Median (ms) & 22.53 & 22.41 & 22.43 & 22.42 & 22.42 & 22.53 & 22.44 & 22.44 \\
Against all runs & 1.004 & 0.999 & 0.999 & 0.999 & 0.999 & 1.004 & 1.000 & 1.000 \\
\bottomrule
\end{tabular}
\end{table}

Three decisions followed. CPU 12 stays in the mask. The measured build is
made on the prepared machine and twice from scratch, tagged only if the
digests agree, because on the unprepared machine a process that is not
pinned can land on the failing core, the compiler included, and a failing
core can compute wrongly without crashing. And tests and analysis on this
machine run after the setup, or pinned away from that core. The last rule
was needed once: after the machine was restored for the hybrid campaign, a
table command stopped with a \code{TypeError} raised inside Python's own
\code{statistics} module on data that the same command reads without error,
a wrong computation with no crash and no kernel log entry. The tables were
written pinned away from the core, and the continuous integration now
regenerates every table on another machine and compares it with the
committed one.

The main campaign adds what a pinned check cannot. With one thread the
kernel put the work on CPU 12 or 14 (Table~\ref{tab:core-split}); the
native filter ran mostly on CPU 12 in 9 of its 30 runs, with a median
1.0003 times that of the configuration. The slow one-thread runs that the
earlier build showed (five filter runs 1.21 times their configuration's
median) did not come back: no native filter run at one thread is above
1.005 times the median.

\begin{table}[htbp]
\caption{One-thread filter runs of the main campaign split by the CPU with the most busy time: runs, their median and slowest run over the configuration's median, and that CPU's share of the run's busy time. CPU 12 is the degraded core.}\label{tab:core-split}
\centering\small
\begin{tabular}{llrrrrr}
\toprule
Environment & Workload & CPU & Runs & Median & Max & Busy share \\
\midrule
Chrome & filter & 12 & 12 & 1.0026 & 1.0144 & 0.643 \\
Chrome & filter & 14 & 18 & 0.9986 & 1.0189 & 0.595 \\
Chrome & filter-seq & 12 & 16 & 1.0043 & 1.0190 & 0.811 \\
Chrome & filter-seq & 14 & 14 & 0.9938 & 1.0216 & 0.807 \\
Firefox & filter & 12 & 11 & 1.0020 & 1.0426 & 0.471 \\
Firefox & filter & 14 & 19 & 0.9997 & 1.0622 & 0.485 \\
Firefox & filter-seq & 12 & 12 & 1.0052 & 1.0525 & 0.562 \\
Firefox & filter-seq & 14 & 18 & 0.9997 & 1.0352 & 0.559 \\
Native & filter & 12 & 9 & 1.0003 & 1.0029 & 0.801 \\
Native & filter & 14 & 21 & 0.9997 & 1.0049 & 0.706 \\
Native & filter-seq & 12 & 22 & 1.0000 & 1.0017 & 0.990 \\
Native & filter-seq & 14 & 8 & 1.0001 & 1.0016 & 0.990 \\
\bottomrule
\end{tabular}
\end{table}

